\documentclass{article}
\usepackage[T1]{fontenc}
\usepackage{lmodern}
\usepackage{graphicx}
\usepackage{amsmath,amssymb}
\usepackage[a4paper,margin=2cm]{geometry}
\usepackage[absolute,overlay]{textpos}
\setlength{\TPHorizModule}{1mm}
\setlength{\TPVertModule}{1mm}
\usepackage[indonesian]{babel}
\usepackage{hyperref}
\setlength{\emergencystretch}{2em}
% unit: Halaman 5

\begin{document}

% === Halaman 5 (llm) ===

\section*{Teknik Analisis Frekuensi}

\begin{itemize}
    \item Pada \textit{cipher} abjad-tunggal, perulangan huruf di dalam plainteks tercermin pula pada perulangan huruf cipherteksnya.
    \item Artinya, huruf yang sering muncul di dalam plainteks, maka huruf cipherteksnya juga sering muncul.
    \item Hubungan statistik antara huruf-huruf di dalam plainteks dengan huruf-huruf di dalam cipherteks menjadi peluang bagi kriptanalis untuk memecahkan cipherteks tanpa mengetahui kunci yang digunakan.
    \item Dengan memanfaatkan tabel frekuensi kemunculan huruf, pasangan huruf (bigram), dan tiga huruf (trigram) di dalam suatu bahasa natural, kriptanalis dapat dapat mencari tabel substitusi huruf plainteks menjadi huruf cipherteks dengan
\end{itemize}

\end{document}
